\section{A width-independent tanh bound without robustness}
\label{sec:productnorm}

A second unrestricted result addresses arbitrary learner width, but with an
exponential dependence on the energy needed by a fixed target circuit.
Its depth dependence is polynomial. The following Rademacher argument
accounts for the fact that tanh is not positively homogeneous.

The normalized-activation and exponential-moment method is closely related
to Nakhleh--Nowak's treatment of deep neural variation spaces
\cite{nakhleh2026}, which explicitly uses $\tanh(st)/s$ and slope
discretization, and to the exponential peeling argument of
Golowich--Rakhlin--Shamir~\cite{golowich2018}. We give a self-contained
product-Frobenius version for the present Gaussian-energy conventions;
we do not claim the underlying proof technique is new.

\begin{theorem}[Product-Frobenius bound for tanh]\label{thm:tanh-product}
Let $\mathcal F_{L,P}$ contain all scalar-output tanh networks with $L-1$
hidden layers, arbitrary finite widths, a linear readout, and product of
layer Frobenius norms at most $P$. Biases are absent except through an
augmented input. On every sample $x_1,\ldots,x_n$ with $\|x_i\|\le R$,
\[
 \widehat{\mathfrak R}_n(\mathcal F_{L,P})
 \le \frac{PR}{\sqrt n}
 \left[2+\sqrt{2(L-1)\log\bigl(3+n(L-1)^2\bigr)}\right].
\]
The bound is independent of every hidden width and of the individual layer
norms apart from their product.
\end{theorem}

\begin{proof}
If $R=0$, every sampled input and every network output on the sample is
zero, so the bound is immediate. Assume $R>0$ below.
Define $\phi_s(u)=\tanh(su)/s$ for $s>0$, $\phi_0(u)=u$, and
$\phi_\infty(u)=0$. These functions are odd, $1$-Lipschitz, and zero at zero.
For finite $s,t\ge0$ and $r\ge0$,
\[
 \phi_s(ru)=r\phi_{sr}(u),\qquad
 |\phi_s(u)-\phi_t(u)|\le u^2|s-t|,
 \qquad |\phi_s(u)|\le1/s\quad(s>0).
\]
The derivative bound follows by differentiating
$\phi_s(u)=u\int_0^1\operatorname{sech}^2(suv)\,dv$.

Let $\mathcal G_0$ contain the identity feature map. Inductively allow
arbitrary finite-dimensional features
\[
 g_j(x)=\phi_{s_j}(a_j^T h(x)),\quad
 h\in\mathcal G_{\ell-1},\quad \sum_j\|a_j\|^2\le1,
 \quad s_j\in[0,\infty].
\]
All these feature vectors have norm at most $R$ on the sample. Put
$B_\ell(\epsilon)=\sup_{g\in\mathcal G_\ell}
 \|\sum_i\epsilon_i g(x_i)\|$ for independent Rademacher signs.
Writing $a_j=r_j u_j$ and using the scaling identity shows
\[
 B_\ell\le
 \sup_{h,\,\|u\|\le1,\,s\in[0,\infty]}
 \left|\sum_i\epsilon_i\phi_s(u^T h(x_i))\right|.
\]
For any $\eta>0$, uniformly on $[-R,R]$ the slope family admits an
$\eta/n$ approximation using at most
$J=3+n^2R^2/\eta^2$ choices: truncate $s>n/\eta$ to the zero function,
and round slopes below that threshold down at mesh $\eta/(nR^2)$.
The error in the Rademacher sum is at most $\eta$.

We use the convex contraction principle
\[
 \mathbb E G\!\left(\sup_{t\in T}\sum_i\epsilon_i\psi_i(t_i)\right)
 \le
 \mathbb E G\!\left(\sup_{t\in T}\sum_i\epsilon_i t_i\right)
\]
for convex nondecreasing $G$ and $1$-Lipschitz $\psi_i$ vanishing at zero.
For completeness, the single-coordinate proof compares
$X_\pm=\sup_t(a_t\pm\psi(t_k))$ and
$Y_\pm=\sup_t(a_t\pm t_k)$. One has
$\max(Y_+,Y_-)\ge\max(X_+,X_-)$ by $|\psi(t)|\le|t|$.
An approximately maximizing pair for $X_+,X_-$ also gives
$Y_++Y_-\ge X_++X_-$, by the Lipschitz property.
Convex monotonicity therefore gives
$G(Y_+)+G(Y_-)\ge G(X_+)+G(X_-)$; iterate over coordinates.
Bounded separable classes follow by approximation from finite classes.
In our application, $u$ ranges symmetrically and $\phi_s$ is odd, so the
absolute supremum is the same as the one-sided supremum.

Apply this principle with $G(t)=e^{\lambda t}$ to each grid slope, then
bound the exponential of a maximum by the sum of the exponentials:
\[
 \mathbb E e^{\lambda B_\ell}
 \le e^{\lambda\eta}J\,\mathbb E e^{\lambda B_{\ell-1}}.
\]
Furthermore, $\mathbb E B_0\le R\sqrt n$, and bounded differences gives
$\mathbb E e^{\lambda B_0}
 \le e^{\lambda R\sqrt n+\lambda^2nR^2/2}$.
Writing $H=L-1$, iteration and Jensen yield
\[
 \mathbb E B_H\le H\eta+R\sqrt n+
             \frac{H\log J}{\lambda}+\frac{\lambda nR^2}{2}.
\]
Take $\eta=R\sqrt n/H$ and optimize $\lambda$ to obtain
$\mathbb E B_H\le R\sqrt n[2+\sqrt{2H\log(3+nH^2)}]$.

Finally, for a particular original network, write
$M_\ell=\|W^\ell\|_F$ and $S_\ell=\prod_{j\le\ell}M_j$.
Dividing its $\ell$th hidden vector by $S_\ell$ changes its recursion to
$\phi_{S_\ell}((W^\ell/M_\ell)g^{\ell-1})$.
Thus its normalized hidden features belong to the enlarged classes above,
and its scalar output is at most $P$ times a norm-one linear readout of
$\mathcal G_H$. Zero-norm layers give the zero function and require no
normalization. This proves the theorem. All suprema are measurable by
countable approximation of finite-dimensional parameter spaces and then
union over widths.
\end{proof}

For the Gaussian energy in our model, the product of the $L$ weighted
squared layer norms is $dP^2$. Consequently
\[
 P\le d^{-1/2}(\kappa_N/L)^{L/2}
 \le d^{-1/2}e^{B/2}\qquad\text{if }\kappa_N\le L+B.
\]
Applying the theorem to the ramp loss
$\ell(yz)=\min\{1,\max\{0,1-2yz\}\}$ gives the following explicit bound.

\begin{corollary}[All widths at bounded excess energy]\label{cor:excess-risk}
Fix $B\ge0$. With probability at least $1-\delta$, every sample-consistent
depth-$L$ network, at any finite width, with training margins at least $1/2$
and energy $\kappa_N\le L+B$ has sign error at most
\[
 \frac{4\sqrt2 e^{B/2}}{\sqrt n}
 \left[2+\sqrt{2(L-1)\log\bigl(3+n(L-1)^2\bigr)}\right]
 +\sqrt{\frac{\log(1/\delta)}{2n}}.
\]
No internal robustness assumption is imposed.
\end{corollary}

\begin{proof}
Augmented Boolean inputs have norm $\sqrt{d+1}$, so $PR\le\sqrt2e^{B/2}$.
The ramp loss is bounded by one and is $2$-Lipschitz. Symmetrization,
scalar contraction, and bounded differences give population loss at most
empirical loss plus four times the network Rademacher bound and the
stated confidence term. Empirical loss is zero and population loss
upper-bounds sign error.
\end{proof}

\begin{lemma}[A scalar transmission construction]\label{lem:scalar-transmission}
For every real $D>0$ and integer $m\ge2$, there are positive scalar weights
$\alpha,b_2,\ldots,b_{m-1},c$ such that
\[
 r_1=\tanh\alpha,\quad r_s=\tanh(b_s r_{s-1}),\quad
 cr_{m-1}=1/2,
\]
and, writing $D'=\max\{D,4\}$,
\[
 D\alpha^2+\sum_{s=2}^{m-1}b_s^2+c^2
 \le(m+12)(D'/4)^{1/m},\qquad c^2=(D'/4)^{1/m}.
\]
\end{lemma}

\begin{proof}
Set $q=(D'/4)^{1/(2m)}$, $k_s=\min\{s,m-s\}$, and
$r_s=q^s/(\sqrt{D'}k_s)$. Then $0<r_s\le1/(2qk_s)\le1/2$ and
$r_{m-1}=1/(2q)$. Choose
$\alpha=\operatorname{atanh}r_1$,
$b_s=\operatorname{atanh}(r_s)/r_{s-1}$ and $c=q$.
For $0<u\le1/2$ the power series gives
$(\operatorname{atanh}(u)/u)^2\le1+u^2$.
Writing $R_s=k_{s-1}/k_s$, the reciprocal pairs
$k/(k+1),(k+1)/k$ satisfy
\[
 \left(\frac{k+1}{k}\right)^2+
 \left(\frac{k}{k+1}\right)^2-2\le4/k^2.
\]
Consequently $\sum_{s=2}^{m-1}R_s^2\le m-2+7$,
$R_s\le2$, and $\sum_{s=1}^{m-1}k_s^{-2}<4$.
Using $q^{2m}=D'/4$ therefore gives
\[
 D\alpha^2\le q^2+1/4,\qquad
 \sum_{s=2}^{m-1}b_s^2\le q^2(m-2+7)+4,\qquad c^2=q^2.
\]
Their sum is at most $(m+12)q^2$. Empty internal sums at $m=2$
satisfy the same estimates.
\end{proof}

\begin{corollary}[Polynomial cost of additional depth in the strong-prior limit]
\label{cor:padded-strongprior}
Suppose a width-$C$, depth-$L_0$ network implements $f$ with global margin
at least $1/2$ and Gaussian energy $E$. Put
\[
 A=\max\{2E,1\},\qquad B=2E+13+7A.
\]
For each learner depth $L\ge L_0$, Corollary~\ref{cor:excess-risk}, with
$B$ replaced by $B+1$, gives a sample bound for the original unrestricted
Gaussian strong-prior learner at every width $N\ge C$.
Specifically, choose $n$ so that the displayed risk bound is at most
$\varepsilon/2$. On one sample event of probability $1-\delta$, for each
$N\ge C$ the posterior Gibbs risk is at most $\varepsilon$ for all
sufficiently small Gaussian scales. The scale threshold may depend on
$N,L,f,\varepsilon$; no uniform threshold is asserted.
The sample bound is polynomial in $L$ and exponential in $E$.
\end{corollary}

\begin{proof}
Balanced cloning into width $N\ge C$ preserves the target function and
increases energy by at most a factor of two. Write $c^2=N\|v\|^2\le2E$ for its readout
cost and $E_{\rm pre}\le2E$ for the other costs.
Apply Lemma~\ref{lem:scalar-transmission} with $D=4c^2$ and $m=t+1$.
Replace the old readout by a hidden layer whose every row is $2\alpha v^T$,
then by the lemma's scalar channels and readout. Its first new preactivation
is $2\alpha z(x)$, whose signed magnitude is at least $\alpha$.
Oddness and monotonicity preserve the sign and final margin. The new first
block has energy $4\alpha^2c^2$, and the remaining costs are exactly the
scalar ones. This gives, for every $t\ge1$, a same-width realization at
depth $L_0+t$ with energy
\[
 E_{\rm pre}+(t+13)\max\{c^2,1\}^{1/(t+1)}.
\]
Since $A^a\le1+a(A-1)$ for $0\le a\le1$ and
$(t+13)/(t+1)\le7$, the target energy is at most $L+B$.
This also covers $t=0$ by cloning directly.

For fixed $N,L$, Gaussian small-variance asymptotics imply that prior
mass on functions of infimum energy greater than $L+B+1$, divided by
prior mass on $f$, tends to zero. There are finitely many Boolean functions,
and $f$ has infimum energy at most $L+B$. This ratio bounds the posterior
mass of those functions for every consistent sample. Relaxed minimum-energy
representatives exist at fixed width, have margin at least $1/2$, and have
the same infimum as strict-margin representatives. Therefore every remaining
consistent function has a representative covered by
Corollary~\ref{cor:excess-risk}. The sample event is common to all widths,
and adding at most $\varepsilon/2$ exceptional posterior mass proves the claim.
\end{proof}

This removes exponential dependence on \emph{added learner depth} for a
fixed target witness, even at arbitrarily large learner widths. It does not
give a bound jointly polynomial in circuit size and depth, and it does not
yet cover unrestricted arbitrary-width posteriors at every nonzero scale.

\subsection{A nontrivial unrestricted family: linear threshold targets}

The preceding result is not limited to padding a fixed network. It gives a
bound polynomial in depth and input dimension for every linear threshold
whose geometric margin on the Boolean cube is inverse polynomial.

\begin{corollary}[Linear thresholds at every learner width]
\label{cor:linear-threshold}
Suppose
\[
 f(x)=\operatorname{sign}(a^T(1,x)),\qquad
 \mu=\min_{x\in\{-1,1\}^d}\frac{|a^T(1,x)|}{\|a\|}>0,
 \qquad D_* =\max\{d/\mu^2,4\}.
\]
For every depth $L\ge2$ and width $N\ge1$, the original unrestricted
Gaussian strong-prior learner admits the following common sample event.
With probability at least $1-\delta$, its Gibbs risk is at most
$\varepsilon$ for all sufficiently small Gaussian scales, provided
\[
 \frac{4 e^{13/2}}{\sqrt n}
 \sqrt{\frac{(d+1)D_*}{4d}}
 \left[2+\sqrt{2(L-1)\log\bigl(3+n(L-1)^2\bigr)}\right]
 +\sqrt{\frac{\log(1/\delta)}{2n}}
 \le\frac\varepsilon2.
\]
The sample event is simultaneous over all widths; the small-scale threshold
may depend on width and depth. In particular the sample bound is polynomial
in $d,L,\mu^{-1},\varepsilon^{-1},\log(1/\delta)$.
\end{corollary}

\begin{proof}
Use the scalar construction with first-layer energy coefficient
$D=d/\mu^2$. If its first scalar weight is $\alpha$, give every first-layer
row the vector $\alpha a/(\mu\|a\|)$. Every first preactivation has the
correct sign and magnitude at least $\alpha$. Put $b_s/N$ in every internal
entry and $c/N$ in every readout entry, where $(\alpha,b,c)$ is the positive
scalar tuple. Oddness and monotonicity preserve the target sign and the
scalar output lower bound. The energy is exactly
$D\alpha^2+\sum_s b_s^2+c^2$, at every width.
Lemma~\ref{lem:scalar-transmission} therefore gives the following upper bound on target infimum energy:
\[
 K_L=(L+12)(D_*/4)^{1/L}.
\]
Functions with infimum energy at most $K_L+1$ have representatives with
product-Frobenius norm at most
\[
 P\le\frac1{\sqrt d}
       \left(\frac{K_L+1}{L}\right)^{L/2}
 \le e^{13/2}\sqrt{\frac{D_*}{4d}}.
\]
Here $(D_*/4)^{1/L}\ge1$ and $(1+13/L)^L\le e^{13}$.
Theorem~\ref{thm:tanh-product}, the ramp-loss argument, and the
small-variance comparison with the target prove the claim exactly as in
Corollary~\ref{cor:padded-strongprior}. The elementary inequality
$\mu\le1$ follows by averaging $(a^T(1,x))^2$ over the cube; hence
$D_*\le4d/\mu^2$, making the asserted polynomial dependence explicit.
\end{proof}

For the finite-variance extension, the same construction has a
\emph{margin-one} center of energy
\[
 K_L^{(1)}\le(L+15)(D_*/4)^{1/L}.
\]
Indeed the scalar construction has readout cost $(D_*/4)^{1/L}$;
doubling that readout adds three times this cost. A posterior cutoff
$K_L^{(1)}+1$ therefore has product bound
$P\le e^8\sqrt{D_*/(4d)}$. For Theorem~\ref{thm:unrestricted-above-one},
\[
 Q_*^2=\left(\frac{K_L^{(1)}+2}{L}\right)^L
 \le e^{17}D_*/4.
\]
Thus that theorem gives a bound jointly polynomial in
$d,L,\mu^{-1},\varepsilon^{-1},\log(1/\delta)$ at \emph{every} learner width
and every exponent $\gamma\ge1+\zeta$, for each fixed $\zeta>0$.
No robustness restriction or infinite-width limit is used. The polynomial
degree may depend on $1/\zeta$; the mean-field exponent $\gamma=1$
remains outside this argument.
